\documentclass[11pt]{article}
\usepackage[margin=1in]{geometry}
\usepackage[T1]{fontenc}
\usepackage{microtype}
\usepackage{hyperref}
\usepackage{xcolor}
\usepackage{enumitem}
\usepackage{booktabs}
\usepackage{tabularx}
\usepackage{ragged2e}

\hypersetup{
  pdftitle={Research Meeting Transcripts (Edited Verbatim)},
  pdfauthor={Nicholas Mino},
  pdfsubject={SURA research meetings with Woody Zhu and Wenbin Zhou},
  colorlinks=false,
  pdfborder={0 0 0}
}

\definecolor{critical}{RGB}{140,20,20}
\definecolor{superseded}{RGB}{120,70,20}
\definecolor{current}{RGB}{20,90,50}
\definecolor{editorial}{RGB}{60,60,60}

\clubpenalty=10000
\widowpenalty=10000
\displaywidowpenalty=10000

\newcommand{\CriticalSource}[1]{%
  \par\noindent\textcolor{critical}{\textbf{EXTREMELY IMPORTANT SOURCE --- #1}}%
  \par\medskip
}
\newcommand{\PlanNote}[2]{%
  \par\noindent\fbox{\parbox{\dimexpr\linewidth-2\fboxsep-2\fboxrule}{%
    \RaggedRight\textcolor{#1}{\textbf{PLAN STATUS (editorial):} #2}%
  }}\par\medskip
}
\newcommand{\Omit}[1]{%
  \par\noindent{\footnotesize\textit{[Omitted: #1]}}\par\medskip
}
\newcommand{\Turn}[2]{%
  \medskip\noindent\textbf{#1 --- #2}\par\nobreak\noindent\ignorespaces
}
\newcommand{\Editorial}[1]{%
  \par\medskip\noindent\textcolor{editorial}{\textbf{Editorial note.} #1}\par\medskip
}
\newcommand{\Takeaway}[1]{%
  \par\medskip
  \noindent\fbox{\parbox{\dimexpr\linewidth-2\fboxsep-2\fboxrule}{%
    \RaggedRight\textcolor{editorial}{\textbf{Editorial summary (not transcript).} #1}%
  }}\par
}

\title{Research Meeting Transcripts\\[0.35em]
\large Edited Verbatim Research Record\\[0.15em]
\normalsize Uncertainty Quantification, Decision Making,\\
and Stability of Optimization-Induced Policies}
\author{Nicholas Mino\\
Carnegie Mellon University\\[0.4em]
\small Faculty mentor: Woody Zhu\\
\small Graduate mentor: Wenbin Zhou}
\date{Document compiled 26 July 2026\\
\small Meetings: 27~Apr, 22~May, 29~Jun 2026}

\begin{document}
\maketitle

\section*{Editorial Policy}
This packet is an archival record of three research meetings. Each meeting is an
\textbf{extremely important primary source} for the research problem.

\paragraph{Retain.}
All substantive research discussion is retained \textbf{edited-verbatim}: wording
follows the source transcript except for the transcription corrections listed
below. Meaning of retained discussion is never altered.

\paragraph{Remove.}
The following are removed and replaced by a specific omission marker:
greetings; Zoom / audio / video troubleshooting; waiting for participants;
screen-sharing and permissions logistics; scheduling unless directly relevant
to research work; SURA/SURF administrative logistics; unrelated conversation.

\paragraph{Speaker labels.}
Device labels in the source (\texttt{Nicholas's iPhone},
\texttt{Nicholas Mino (Posematic)}) are normalized to \textbf{Nick Mino}.
Faculty is labeled \textbf{Woody Zhu}; graduate mentor \textbf{Wenbin Zhou}.

\paragraph{Transcription corrections (only).}
Obvious ASR errors are corrected without changing meaning. Documented
corrections include: Chu$\to$Zhu; Minow$\to$Mino; Shura/Sur$\to$SURA;
iTriton$\to$Triton; go part$\to$grid part; iClear$\to$ICLR;
confrodiction$\to$conformal prediction; brief skin$\to$brief skim;
overlapping draft$\to$Overleaf draft; Women$\to$Wenbin; JetGPT/JGPT$\to$ChatGPT; AIS system$\to$AI system;
Tiana$\to$Angelopoulos; author message model$\to$optimization model;
height of intuition$\to$high-level intuition;
automatic models$\to$optimization models (in context);
pre-paper$\to$CREAM paper (in context). Light punctuation and
filler cleanup that does not change meaning is allowed. Ellipsis
(\ldots) marks intra-turn cuts of non-research fragments only.

\paragraph{Editorial apparatus.}
``PLAN STATUS,'' ``Editorial note,'' and ``Editorial summary'' boxes are
\textbf{not} transcript text. They exist to prevent superseded plans from being
mistaken for the current agenda.

\section*{Sources (all critical)}
\begin{enumerate}[leftmargin=*,itemsep=0.35em]
  \item \textbf{Meeting~1} --- SURA: nmino + Dr.~Zhu \quad Mon, 27 Apr 2026\\
        Source: \texttt{SURA\_ nmino + Dr\_ Zhu Transcript(1)}
  \item \textbf{Meeting~2} --- Nick / Wenbin \quad Fri, 22 May 2026\\
        Source: \texttt{Nick \_ Wenbin Transcript(1)}
  \item \textbf{Meeting~3} --- Nick / Wenbin / Woody \quad Mon, 29 Jun 2026\\
        Source: \texttt{Nick \_ Wenbin \_ Woody Transcript}
\end{enumerate}
Also consulted during compilation: a prior rewritten draft and the earlier
transcript packet PDF.

\section*{Plan Evolution (Read First)}
Later meetings supersede earlier proposals. Do not treat early proposals as the
current agenda.

\begin{center}
\begin{tabularx}{\linewidth}{@{}l >{\RaggedRight\arraybackslash}X l@{}}
\toprule
\textbf{When} & \textbf{Idea / direction} & \textbf{Status} \\
\midrule
Apr 2026 & Lab framed by three pillars: (1)~UQ for high-stakes decisions;
  (2)~grid-resilience applications; (3)~human--AI collaboration.
  Nick to read papers and choose a pillar. &
  \textcolor{current}{\textbf{Still foundational}} \\
\addlinespace
May 2026 & Nick chooses \textbf{algorithmic} pillar (CREDO / CREAM). &
  \textcolor{current}{\textbf{Still foundational}} \\
\addlinespace
May 2026 & Proposal~A: quantization $\to$ risk-profile comparison
  (engineering / testing). Woody redirected toward pure algorithmic decision
  analysis. &
  \textcolor{superseded}{\textbf{NO LONGER THE PLAN}} \\
\addlinespace
May 2026 & Proposal~B: \textbf{adaptive regret / certified bounded-regret
  frontier}. Overleaf formulation assigned. &
  \textcolor{superseded}{\textbf{NO LONGER THE PLAN}} \\
\addlinespace
May 2026 & Wenbin backups: (i)~CREAM in sequential / non-i.i.d.\ setting;
  (ii)~accommodate model / solver error. &
  \textcolor{superseded}{\textbf{Backup only; not active}} \\
\addlinespace
Jun 2026 & \textbf{Current direction:} stability of policies induced by
  optimization; stabilized selection; formal problem setup; AI predictors +
  downstream optimization; ``should I trust this decision?'' as hypothesis
  testing. Priority reading: prediction-powered inference; algorithmic /
  post-selection stability (Angelopoulos \& Jordan line). &
  \textcolor{current}{\textbf{CURRENT PLAN}} \\
\bottomrule
\end{tabularx}
\end{center}

\clearpage
\section{Meeting 1 --- Lab Framing and Three Pillars}
\textbf{Date:} Mon, 27 Apr 2026 \quad
\textbf{Participants:} Nick Mino, Woody Zhu\\
\textbf{Source title:} SURA: nmino + Dr.~Zhu

\CriticalSource{Meeting 1 (SURA kickoff with Dr.~Zhu)}
\PlanNote{current}{Introduces the lab's three research pillars and the
expectation that Nick will read publications, choose a pillar/methodology,
and (ideally) propose an unexplored follow-on idea. No concrete project title
yet---this meeting sets the problem space, not the final formulation.}

\Omit{0:01--2:16 audio checks, AI note-taker remarks, waiting}

\Turn{2:17}{Nick Mino}
Just in case Professor Zhu isn't able to come to the meeting, but might read
this AI meeting notes transcript. I'm Nick Mino. I'm a sophomore in AI and SCS.
I'm minoring in algorithms and complexity. And I'm very interested in this
project because I've had a lot of experience or a good amount of experience in
like machine learning implementation. But I think the thing that I've learned
from that is that I want to go more in the research route.

\Turn{2:57}{Nick Mino}
I have experience in PyTorch and deployment through Triton servers and some
Docker and Linux experience as well. I have strong mathematical foundations
from 15-151 and I'm in 36-219 probability theory. Yep. I'm very interested in
this project as well, because from what I understand, it's about uncertainty
modeling and decision making as it relates to energy infrastructure. From
looking through some of Professor Zhu's research, he's done work on things like
power grids. And I live in California. I have PG\&E. They could definitely use
some help with that kind of work. And I want to do work that helps my community
in looking into more mathematical grounded machine learning, which this is
compared to what I've been doing so far.

\Turn{4:14}{Nick Mino}
And yeah, my strengths are in implementation, but I do want to closer to the
math theory and theory in general.

\Turn{4:27}{Nick Mino}
I can devote a lot of time and I'll be on campus.

\Turn{4:59}{Nick Mino}
Another reason I'm interested in doing research this summer, specifically, is
that after doing all this implementation, I think that the route I want to go
is into research. So I'm focusing on research now. That I can hopefully get a
PhD.

\Omit{5:33--6:07 greeting after Woody joins; late-arrival / meeting-link remark}

\Turn{6:08}{Nick Mino}
Yeah, I'm a sophomore in artificial intelligence. I'm minoring in algorithms
and complexity. And for the summer, I'm looking to do research specifically. I
have a lot of implementation experience. I've worked as an ML engineer at a
startup.

\Turn{6:26}{Nick Mino}
And after doing that, I think that I've learned that while I appreciate the
implementation part, I like the research part more. I've talked to some
research engineers from DeepMind and their experiences told me has made me
think that I like research more. I've done some research as well. I did some
paleobiology research at Stanford by data modeling. Okay.

\Turn{6:53}{Woody Zhu}
And what type of things that you like or interest you the most, just out of
curiosity? I know you're interested in ML, math, maybe stats as well, but for
example, what's What particular methodologies that interest you the most or
what specific applications that you feel most interested in?

\Turn{7:16}{Nick Mino}
I'm pretty eco-conscious. I like projects that relate to green energy and
things. I did a hackathon where we won the most sustainable project for this
thing that decreased LLM token usage. Like compacting chat history in a
specific way and routing models. Okay. And I did the paleobiology research
because, you know, I care about the environment. That was very interesting to
learn about fossil displacement.

\Turn{7:51}{Woody Zhu}
Okay. What type of course that you have taken so far with respect to, let's
say, machine learning, statistics, optimization and the math?

\Turn{8:02}{Nick Mino}
I took discrete mathematics 15-151. Okay. 21-241 is our linear algebra. Okay.
Right now I'm in 36-219 probability theory and it covers basic probability
theory up to random processes.

\Turn{8:21}{Woody Zhu}
Over the summer I'll be taking statistics and 3D calc Got it. On campus.

\Editorial{ASR attribution on 8:21 is ambiguous in the source (Woody labeled;
content is Nick stating summer courses). Sense of the exchange: Nick will take
statistics and multivariable calculus on campus over the summer.}

\Turn{8:30}{Woody Zhu}
So you'll be staying around the campus over the summer if I understand
correctly.

\Turn{8:35}{Nick Mino}
Yeah. Okay.

\Turn{8:37}{Woody Zhu}
So yeah, so if you will be remaining physically here on campus, then I can give
a thought and trying to come up with a project that fits your interest the
best. But before we go into that, so you're looking for SURA, right? Mentorship
through the SURA program, if I'm not wrong. Yeah.

\Omit{9:01--9:09 SURA application deadline logistics}

\Turn{9:10}{Woody Zhu}
So what I suggested to do is to go ahead and read the papers that we have in
our publication list. So you can just take a look at those papers that is
either published recently or the paper that is current still ongoing. So we
have a list of the preprints and ongoing papers, but I also have some recent
papers been published in 2026. So you can just take a look at this papers and
to see if you feel comfortable with the depth of the methodology and the
theory, because usually a student feel very enthusiastic about the research,
but whenever they really get their hands on the project, they feel like, oh,
This is not what I'm looking for. I'm not expecting that much as an undergrad
research project. So what I suggest to do is to look at those papers and to get
the sense what type of methodologies we might need to use in our research
project, and also to get the sense what type of applications we're dealing
with. And then you can get back to me and say, hey, Woody, I read the paper
XYZ. I feel like, well, this aligns my interests really well. And I can give
you some ideas about maybe the title and abstract for SURA project for your
application based on your interest, and then we can start things from there, if
that makes sense to you.

\Omit{10:36--10:45 brief exchange on locating publications on Woody's webpage}

\Turn{10:46}{Nick Mino}
Yeah, I think I've been looking around a bit on there before. You've done some
research on uncertainty modeling for energy grids, which is why I brought up
the grid part.

\Turn{10:58}{Woody Zhu}
Yeah, exactly. So again, I think our research consists of three major pillars.
So one pillar is about uncertainty quantification for high stake decision
making, right? And mostly motivated by the application from energy systems. So
the high level idea is we want to figure out a way to tell people how uncertain
is the future in a way that can lead to better decision or most robust
decisions. So one of the good example I can give you is just put yourself in
the shoes of home buyers, right? You want to make a purchase of the property
and oftentimes you have to mentally you need to predict how much is the
property, right? You have to make a prediction about the price of the house and
then based on the prediction you will make a decision I say, hey, if the house
is worth, I don't know, \$500{,}000, I happen to have \$500{,}000 in my pocket,
I'm going to go ahead and make the purchase. So that's just the way how we
usually make our decision. And if this is like a high stake decision, meaning
I'm kind of risk averse, it usually implies that during that particular
prediction phase, I have to be a little bit more conservative. Instead of
making a prediction of the price of the property, I'm going to anticipate what
is the worst case scenario, like how expensive this house could be in the worst
case, right? So usually the way how we approach it is to construct a confidence
interval, right? For example, the house worth, I mean, the true price for this
property may be like \$500{,}000, but what I'm going to do is to construct an
upper bound and lower bound. Maybe the upper bound could be \$600{,}000 and
lower bound could be \$400{,}000. And then you're going to take actions in
response to the worst case scenario, which is \$600{,}000. If it is above your
budget, then we're going to say, hey, I don't want to buy it because there is a
chance that it's going to go above my budget. If your worst case is lower than
\$500{,}000, you can take an action more confidently, because you know, at
worst case, you can still afford it. So here really leads to our research
question. How we can construct this uncertainty band in a way that can give us
the best result? So it involves a good understanding on the optimization, and
more specifically, robust optimization. But most importantly, it involves a
good understanding on the statistical method and the machine learning models,
because you This is something goes beyond the traditional ML method because in
ML, people are usually less concerned about uncertainty, more concerned about
predicted accuracy. If you're familiar with ML models, right, they usually just
minimize the predicted error. But now we want to go beyond the minimization of
the prediction error. Instead, we want to focus on more on how confident we can
quantify the uncertainty that is given by a predicted model like, you know,
logistic regression or linear regression.

\Turn{14:15}{Woody Zhu}
So this is one pillar. The other pillar is more like application driven. We're
working very closely with the utility and regulators that's trying to help them
build a model to answer the practical questions with respect to the enhancement
of the grid resilience. So that's kind of the major focus of our application
domain. And last but not least, we also care a lot about human-AI collaboration
in high-stakes decision making, because in the future, we can anticipate the
way how we make our decisions is not completely dependent on ourselves, but
also dependent on how we utilize the knowledge and expertise coming from AI
agents. So how we can come up with a paradigm that can enable efficient and
effective collaboration between human and AI agents is something we also are
interested in. So there are the three line research agenda, as you can tell
from our publication, and you can just look at these papers that falls within
each one of the category, and you can decide which one interests you the most.
And then you can get back to me and tell me, you know, which one you like
better than the others. And then we can start it from there. I can give you
more concrete ideas about the title, the abstract of the program, about your
project and then you can make use of those information to submit your
application if that makes sense.

\Omit{14:15 continued: weekly/biweekly meeting cadence; group meetings;
visitor/undergrad cohort logistics; SURF/CERF funding pipeline}

\Turn{17:08}{Nick Mino}
I think I have one or two questions here. I wrote them down. Go ahead.

\Turn{17:18}{Nick Mino}
What does a strong SURA outcome look like, hope students to accomplish by the
end of the summer?

\Turn{17:27}{Woody Zhu}
I have to be honest with you, I have no expectation, especially for students at
second or even first year. I think the main purpose is trying to inspire the
students who are interested in the research. So it's more like a long-term
investment for myself. I wanted to invest at least like a tiny portion of my
time over the summer to the students who potentially could be a leading
researcher in the future. And I wanted to see if they're interested. If they're
interested, definitely I would like to spend more time with the students maybe
in the next summer through the CERF. And I have very luckily advised a few
other students who have successfully, you know, get the grants from the CERF
and also because of the CERF, they get a really nice PhD offers from other
universities.

\Turn{18:21}{Woody Zhu}
So yeah, that's the things I'm expecting. So I do not expect anything
concretely from the students on the research project. Instead, I'm more
interested in seeing whether or not you're interested in the research.

\Turn{18:35}{Woody Zhu}
As long as you feel, oh, I really love research. I enjoyed working with Woody
over the first summer. That's what I'm expecting. But if you don't like it, let
me know, no pressure at all. Um, because if you don't love something doing that
kind of stuff, just waste of a time. Okay.

\Turn{19:00}{Nick Mino}
And I think I had one more. Yeah. Um, how independent would we be? You say we
meet weekly or biweekly. Would we be doing like our own mini project or would
we be working? That's a great question.

\Editorial{Trailing ``That's a great question.'' in the source is likely Woody;
retained attached to Nick's turn as in the raw transcript.}

\Turn{19:17}{Woody Zhu}
So the way how it works within our group is I'm normally going to give students
a mini project so the student can work independently on this project. Because
by the end of the day, what really defines a researcher is your independence,
right? If you can independently come up with the idea and figure out the
solution and turn it into academic publication, that would be the entire
training pipeline you have to go through as a researcher in the future, no
matter which discipline you're looking at, right? So this is a kind of
standard. But on top of that, I'm also encouraging students to participate in
group meetings, and other projects, you don't necessarily have to do anything.
You can just be the observer to see how things are moved forward by the other
senior students. So you can learn from the interaction with those senior
cohorts. But if you have the capability to make your own contribution, that
will be highly recommended and highly welcome, but it's not required. It varies
from person to person. Some people just get things very slowly at at the
beginning, then all of a sudden, once they figure things out, they can be super
productive. Or some people, they're more enjoying the teamwork rather than
working on a single project themselves. So I always change the strategies along
the way. And the strategy may vary from student to student.

\Turn{20:54}{Nick Mino}
Thank you. Yeah, those are the questions I had. Sure.

\Turn{21:02}{Nick Mino}
I'm going to read through your papers and come back to you with the things that
I found most interesting, like methodologies and which of the three pillars.

\Turn{21:13}{Woody Zhu}
Right. Just keep those three pillars in mind and read those papers accordingly
and get things back to me about what are you interested in, if you're
comfortable with the methodologies, with respect to to the pillar that you're
interested in. And also, it would be even better if you can come up with some
of your own ideas to let me know, hey, Woody, I read your paper. I feel like
there's one exciting area that is unexplored by your paper. Can we work on that
together? That would be awesome, but not required again. OK? Yep. That sounds
great. Thank you. You have a very tight timeline, so I would expect that you
have to things back to me within a week. Otherwise, you wouldn't be able to
make up the deadline. But again, it's up to you. No pressure at all.

\Omit{21:13 continued closing pleasantries / goodbye (22:16)}

\Takeaway{Three pillars define the lab. Core research question on pillar~1:
construct uncertainty bands that support robust high-stakes decisions (beyond
minimizing prediction error). Nick to select a pillar and, ideally, propose an
unexplored idea.}

\clearpage
\section{Meeting 2 --- Algorithmic Pillar; CREDO/CREAM; Adaptive Regret Proposal}
\textbf{Date:} Fri, 22 May 2026 \quad
\textbf{Participants:} Nick Mino, Wenbin Zhou\\
\textbf{Source title:} Nick / Wenbin

\CriticalSource{Meeting 2 (Nick / Wenbin project scoping)}
\PlanNote{superseded}{At this meeting, the \textbf{active plan} was Nick's
adaptive / certified bounded-regret frontier idea, to be written as a
mathematical problem formulation in Overleaf (CREAM notation as starting
point). Wenbin also offered sequential-CREAM and error-accommodation ideas as
\textbf{backups}.\\[0.35em]
\textbf{IMPORTANT:} As of Meeting~3 (29 Jun 2026), the adaptive-regret-frontier
proposal is \textcolor{superseded}{\textbf{NO LONGER THE PLAN}}. Retain this
meeting for intellectual history, CREDO/CREAM context, rejected directions
(quantization-as-engineering), and backup ideas---not as the current agenda.}

\Omit{0:03--1:25 greetings; summer-class small talk; SURA vs SURF credit/pay
logistics}

\Turn{1:26}{Wenbin Zhou}
\ldots I think the most important thing for research is about like interest, is
the interest to like dig into like one thing to. Go really deep into that and
don't get frustrated when you meet like problems, especially in like this field
of like in this era of like AI when you have like these LLMs to help you. I
don't really think like technical knowledge is that important anymore because
you can literally just prompt LLMs to teach you how to write code, teach you how
to do the math and teach you a lot of things. Things. So yeah, I think
motivation is like the key thing. And I'm glad you're very motivated. Thank you.
So yeah, so could you tell me a little bit about what you have and what you have
discussed in your like previous meetings and what kind of like papers you have
already read? I just want to know like some sort of like background knowledge on
what what we're expecting for for this sort of like collaboration together.

\Omit{2:40 opening: SURA hour-commitment logistics (minimum 90 hours)}

\Turn{2:40}{Nick Mino}
\ldots first our meeting was more of a, he gave me some papers to read from was
the like AI human decision making like collaboration process was the first
pillar. The second pillar was like the algorithmic part. And then the other
pillar was its applications to the spatial temporal models like energy grids?

\Turn{3:44}{Nick Mino}
So I've been reading\ldots\ A paper.

\Turn{3:49}{Wenbin Zhou}
Could you tell me the exact name of the paper, or do you remember who wrote it?
Yeah.

\Turn{3:55}{Nick Mino}
Well, I ended up choosing the algorithmic one. I read the CREDO paper. Okay.

\Turn{4:03}{Wenbin Zhou}
CREAM paper.

\Turn{4:04}{Nick Mino}
Well, I read, I'm still, you know, uh, trying to to go more into the science
because it's it's not my area of expertise yet. But um, yeah, um, I'm also, I
mean, I recently, I think I have a list here.

\Turn{4:26}{Nick Mino}
got some more papers that I plan on reading that I looked around apparently
good for like background information.

\Turn{4:37}{Nick Mino}
Yeah, yeah. So, the specific papers were CREDO and CREAM.

\Turn{4:42}{Wenbin Zhou}
Yeah, those are the paper that I wrote. So, that's good.

\Turn{4:55}{Wenbin Zhou}
Like after the papers, was there any like specific thing Woody said he wanted
to do or anything?

\Omit{4:47--4:52 brief acknowledgment that Nick saw Wenbin's name on the papers}

\Turn{5:03}{Nick Mino}
He suggested that I well, he said that he would like it if I like came up with
an idea that wasn't previously explored in the papers. And I was interested in
some of my engineering work, I've messed around with like quantizing models.
And I was interested in how different quantization techniques might not affect
accuracy very much in comparison to each other, but might affect a risk
profile, because that seems like an interesting, important area to study. So
that was something I suggested to him. And then his feedback was. To make it, I
think I have, I don't know what exactly in the email, but it was to make it
more closely like the pure algorithmic like decision analysis part rather than
as much like engineering type applied stuff, which is what my original proposal
was, because mostly like testing and comparing.

\Turn{6:14}{Nick Mino}
But I want to do the algorithmic research.

\PlanNote{superseded}{Quantization $\to$ risk-profile comparison (engineering /
testing framing) was Nick's first proposal and was redirected by Woody.
\textcolor{superseded}{\textbf{NO LONGER THE PLAN}} as a primary project
(later partially reappears only as a special case of ``model error'' in
Wenbin's backup ideas).}

\Editorial{Source labels 6:17 as Wenbin Zhou, but the wording (``my second
proposal'') continues Nick's project narrative. Attribution corrected to Nick.}

\Turn{6:17}{Nick Mino}
So we ended up, my second proposal was to make an adaptive regret risk frontier
for like bounded, certified bounded regret. So say a model, you know, is a
99\% chance.

\Turn{6:37}{Nick Mino}
or it's a 98\% chance that the regret is going to be maximum of like 100 units,
but then there's a 99\% chance it's going to be under 10{,}000 units. Well,
knowing that very big jump would be useful. So, if I could make an algorithm
that basically creates like an adaptive frontier, so it sorts out where the
biggest like jumps are, and then gives bounded regret for those that might be
useful for. Decision making because in a lot of the decision making stuff, I
saw mostly focus on average regret or expected regret.

\PlanNote{superseded}{Adaptive / certified bounded-regret frontier was the
\textbf{active plan after Meeting~2}. Overleaf formulation was the next step.
As of Meeting~3, this direction is
\textcolor{superseded}{\textbf{NO LONGER THE PLAN}}.}

\Turn{7:16}{Wenbin Zhou}
Yeah, I see. That's a good idea. That's a good direction to pursue. Yeah, I
think that sounds good. Also, regarding, so I think what you have just
mentioned is actually very closely related to like a cream, the cream paper,
right? So, with the how you want to construct like this sort of regret profile,
I would expect, like, also, if you were to achieve that sort of task, you want
to construct this sort of like Pareto frontier and also to sort of see how this
can be evaluated. I think this are, these are like, yeah, I think this is good
because both Credo and Cream. So Credo was accepted by ICLR and also Cream was
accepted by ICML. So these are like good ML conferences. So yeah, I think we
definitely can do some like follow-up works on these. I think they all have
great potentials. Yeah, and also aligns with your with your background. Like, I
think you're doing like a CS major, right? In or like in AI, yeah\ldots\ And
aside from this idea that you brought up, I also actually personally have some
other ideas regarding some kind of follow-up work, especially regarding like
cream. These were like
some things that when I was doing like the rebuttal, so like when we do like
sort of reviewers review my paper, they kind of challenge some of the parts of
the paper. I have to give them like a rebuttal, which is to clarify why we
think like these challenges can be resolved or like it's already resolved in
the paper. So, there were some interesting ideas that they had mentioned. So,
like, one thing is that maybe like cream can be extended to sort of like a
dynamic adaptive setting where if your decision can, if you're sort of like is
making like a sequential decision. So, um, the current setting here assumes
everything to be like IID. So, like, you making decisions independently, but in
many real-world applications, your decisions must be sequential. So, then, how
can we ensure that our constructed like Pareto Frontier to be valid in this
sort of sequential setting? This is a good idea, in my opinion, especially for
you, because I think this is. Kind of like straightforward. This problem is
very well formulated because, yeah, as I've just said, this is basically just
extending this cream framework into the sequential setting. And there are
already some useful literatures on the sequential side that's something we can
leverage. And yeah, so. The other idea is that can we consider like some sort
of like error, maybe like errors in different components? This kind of relates
to like your quantized model idea, right? So like when you, so like just to
clarify, when you mentioned like quantized, you're just basically saying like
when your model is able to only like output like numbers in some sort of
resolutions. Is that what quantitize means? I'm not super familiar with that.

\Turn{11:07}{Nick Mino}
It's when you take the weights and say they're in like 16-bit, you make them
8-bit. So you have less accuracy, but you still hope.

\Turn{11:18}{Wenbin Zhou}
Yeah. So like approximation. Yeah, yeah. Okay. Well, yeah, I think then it's
the same idea. So basically, it means your model is. Is not accurate due to
some like real-world physical constraint because of like this sort of
quantization. Yeah, so basically, this problem can your idea of quantitization
can be like generalized to like sort of considering if there is some kind of
error in this kind of model that makes it not accurate, how can we really
adjust for this sort of like estimation? This is also like a generic problem
that. Framework can look into, how to accommodate different errors. One, aside
from your quantitization model idea, another thing that we can look into is
that, for example, when your optimization solver is not accurate, is there any
way we can adjust this frontier in any ways? So, the reason why there's, I
mentioned optimization solver, is because, like, if you read the CREAM paper,
so there. Is one step where you have to solve for like the regret. So you have to
solve for the optimization problem and get that number. And there might be some
errors without solvers because they might not always give you the exact
solutions. So how can we approximate that frontier during that setting? This is
also one possible idea.

\PlanNote{superseded}{Wenbin's two ideas (sequential CREAM; model/solver error)
were explicitly framed as \textbf{backup plans} if Nick's adaptive-regret idea
did not work out. They were not adopted as the Meeting~3 direction either.}

\Turn{12:50}{Nick Mino}
Yeah, I'm totally open to switching project ideas to whatever you think is
best. Yeah, I'm. Yeah, yeah.

\Turn{13:04}{Wenbin Zhou}
Yeah, no pressure on that. So, yeah, what I was just saying are like, I guess,
some backup plans. These are things because. Like, I know it's challenging for
like undergrads to work on like projects, especially if you want to think
independently of an idea and try to work on it. This might be challenging, but
it's also rewarding because this gets you to taste like the raw feeling of
research. This is like essentially how an independent scholar should be doing.
But in case you need more guidance, what I just mentioned are two
well-formulated ideas that I'm personally confident about doing. And I'm
confident I can work with you. I can guide you in writing out something like a
successful manuscript or even turning into some sort of submission, like
conference submission. If you want a publication during your undergrad, which
is. Like, I would think a good thing if you consider to apply for like grad
schools or jobs.

\Turn{14:10}{Wenbin Zhou}
So, yeah, but yeah, I think maybe as for the next set plan, I think it's a
since I think you already have like some sort of proposal on your side about
this adapted, I adapted idea. So I I would suggest like as a next step, maybe
we can try to. I'm not sure if you used, do you know how to use LaTeX? Yes.

\Turn{14:39}{Nick Mino}
So, I think maybe as a next step, we can open like an Overleaf draft together.

\Editorial{Source labels 14:39 as Nick, but the Overleaf suggestion continues
Wenbin's next-step plan; Nick's audible reply is assent. Sense preserved.}

\Turn{14:47}{Wenbin Zhou}
And maybe you can put down some of the formulation of that idea into that
draft. And by formulation, what I mean is to. Try to write out everything
mathematically. A good practice is to like first start off by referencing all
some of the notations that is already used in the cream paper, like how I set
it up the problem. You can just sort of like copy the sort of like text and
also like the notations here and try to make some kind of incremental changes
to express what you. Have proposed. Like, what do you mean by sort of adaptive
choice of regret in there? So, once we have that sort of clear formulation of
the idea, it will be more easier for me and also Woody to help you to think
about what we can do next. Right?

\Turn{15:48}{Nick Mino}
Yeah, I think that I think that's a cool next step here.

\Turn{15:52}{Wenbin Zhou}
And we can see if this idea eventually works out in case if it doesn't, I think
it would be always possible for us to sort of fall back to these sort of ideas
that I mentioned. Yeah.

\Omit{16:13--16:49 whether Woody had other plans; Woody travel / return timing}

\Turn{16:50}{Wenbin Zhou}
So for the next step, I think maybe you can open like an Overleaf draft\ldots\
And as a first step, you can start with like a section that's called like a
problem formulation. And try to write out like things mathematically, how you
want to set up the problem that you raised in your proposal. And yeah, maybe we
can aim for like after you finish it, we can then schedule like another short,
quick chat meeting. I can help you look over the formulation and yeah, and see
we can finalize that. Before we meet with Woody. So, like, by the next time we
meet with Woody, we can present to him like this sort of formulation and get
some feedback from his end.

\Turn{17:42}{Nick Mino}
Yeah, that sounds good.

\Turn{17:48}{Nick Mino}
Yeah, I mean, I have like one or two questions about. Okay. So, you know,
thinking about my problem. I'm not like, I'm not very experienced in the field,
so I'm not 100\% sure how useful my like suggestion would be in practice. Cause
I really want to do research, but I also want to like, you know, do impactful
research that isn't just like a random idea I had that's novel, but not, you
know.

\Turn{18:36}{Wenbin Zhou}
Yeah, so could you elaborate on that or what you're trying to express?

\Turn{18:43}{Nick Mino}
Yeah, so I just wanted to know if like the adaptive like regret bound frontier
idea is would like would it actually be something useful in decision making, or
am I being like naive about the usefulness of regret bounds changing sharply?

\Turn{19:13}{Wenbin Zhou}
Well, yeah. Well, first of all, I think this is a good question that you
raised. And this is one thing that. Should be aware of: is that research itself
is actually full of uncertainty. So, like, when you have an idea, it's normal
to think about how, like, whether or not this idea is going to work out,
whether it's impactful, whether it's going to do this or do that, whether it's
going to make us successful. So, these are all the things that you can think
about before you write things out. Usually this is a bad thing if you become
hesitant because of these things, because you haven't even implemented it or
like wrote it out. So, there's no way to even judge whether or not this idea to
work out. So, my suggestion is always to encourage you to start doing things
first to just finish doing things and then look back at whether or not it works
well, right? Don't let these kind of like difficulties or challenges that you
think might be ahead of to impede like your impede your process of like doing
things like to harm your productivity. And yeah, this is the first thing. And
the second thing is that. Actually, like, just by plain language, I wasn't able
to fully understand what you actually mean by adaptive. So, this is also why I
want you to write out this formulation so I can give it a closer look. Like,
whether or not this thing will be impactful, it is me and also Woody, we can
both look at it and decide. And also, this is not, it doesn't, it's not too
costly, right? I think writing out the formulation would. At most take you like
one week, so this is not a too time-consuming thing to do. So, even if it
doesn't work out, um, we can then, if it doesn't work out, we can just stop
pursuing this direction, we can change our formulation, perhaps to like
something that I mentioned, which is somewhat guaranteed to work, yeah, but um,
yeah. So but I would still encourage you to think independently on that process
because I think essentially what you came for for the SURA project is to get
some sort of like training on research. And I and also Woody, I don't think to
deprive this sort of like interest, the most interesting part of research,
which is to think about what kind of problem you want to work on from you. So
instead of like just assigning you some sort of project we want to do. At least
we want you to like first try, try out to sort of like think about this project
problem yourself, and then to use whatever we think as like a backup thing,
right? So, at least this is what I'm planning.

\Turn{22:13}{Nick Mino}
Yep, that sounds great. Thank you.

\Omit{22:15--22:31 paper-sharing logistics while Wenbin searches}

\Turn{22:32}{Wenbin Zhou}
Because I'm sure you're not too familiar with. So, first of all, conformal
prediction, that's not something you're familiar with, right? No.

\Turn{22:41}{Nick Mino}
Okay.

\Turn{22:42}{Wenbin Zhou}
And familiar with like optimization, especially in the context of like
operations research? Not the ones you see in like at machine learning, like
SGD, not these kind of optimization, but like linear programs and stochastic
programs.

\Turn{23:01}{Nick Mino}
Oh, I don't have experience in that.

\Omit{23:08--23:50 Wenbin searching for tutorial papers; linear-program tutorial
hand-off logistics}

\Turn{23:51}{Wenbin Zhou}
So I think in the meantime, while you are doing that sort of like formulation
for this problem that you want to solve, I will also recommend to like read
these sort of like paper to get some like background knowledge. I expect like
these are things that we might need to use like down the line. Down the road
for this research. So, one is conformal prediction, and the other is like some
kind of optimization. Okay, this seems like an okay paper.

\Turn{24:31}{Wenbin Zhou}
Yeah, okay, so yeah, please save these two paper for like your future
reference. You can like look at them like while you're writing the problem
formulation. There's no pressure, you don't have to necessarily finish it in
one week. And you don't have to go into too much details of this paper. Just
like going through it with a brief skim will be okay. And also, these papers
shouldn't be too relevant with the problem formulation that you want to write.
So you don't have to be like, so it doesn't like have to go sequentially, like,
you have to first read the paper and then write the problem formulation. No,
you can write the formulation even without knowing the things in these paper,
yeah.

\Omit{25:20--26:23 Overleaf invite logistics (send to Wenbin first while Woody
traveling); closing}

\Takeaway{Algorithmic pillar selected (CREDO/CREAM). Active proposal at the
time: adaptive certified regret frontier + Overleaf formulation. Background
reading: conformal prediction + OR-style optimization. Backups: sequential
CREAM; accommodate prediction/solver error. \textbf{Status now:} those plans
are \textcolor{superseded}{\textbf{not the current plan}} after Meeting~3.}

\clearpage
\section{Meeting 3 --- Stability of Optimization-Induced Policies}
\textbf{Date:} Mon, 29 Jun 2026 \quad
\textbf{Participants:} Nick Mino, Wenbin Zhou, Woody Zhu\\
\textbf{Source title:} Nick / Wenbin / Woody

\CriticalSource{Meeting 3 (Nick / Wenbin / Woody --- current research direction)}
\PlanNote{current}{This meeting establishes the \textbf{CURRENT PLAN}. The
research focus is no longer the adaptive regret frontier. It is:\\[0.25em]
(1) formally define \textbf{stability of a policy induced by an optimization
problem};\\
(2) find a \textbf{stabilization mechanism} enabling valid optimization-based
policies / post-selection-style inference;\\
(3) application vision: AI (e.g.\ LLM) predictions feeding a downstream
optimizer; end product framed as hypothesis testing---``should I trust this
decision?'';\\
(4) priority reading: prediction-powered inference; algorithmic /
post-selection stability (Angelopoulos \& Jordan); continue literature survey
of stabilized selection methods.\\[0.25em]
Earlier proposals (quantization engineering; adaptive regret frontier;
sequential CREAM backup; solver-error backup) are
\textcolor{superseded}{\textbf{NO LONGER THE ACTIVE PLAN}} unless later
explicitly revived.}

\Omit{0:05--1:45 screen-sharing permissions / Zoom host troubleshooting;
assignment paperwork mention}

\Turn{1:46}{Nick Mino}
Yeah. So from working on the algorithmic stability. Framework. I kind of
decided on the notation for a stabilized selection rule to look like this
because I couldn't really find any specific examples. So I consulted some AIs
and they said this was the closest to standard that I could find.

\Omit{2:19 waiting while Woody reads the shared notation}

\Turn{2:19}{Woody Zhu}
They denote the observation data set and case on decision A would be a data
dependent selection strategy and as hats denote the selected object.

\Turn{2:45}{Woody Zhu}
We seek a stabilized selection rule, a theoda producing okay where a theta is
the stabilized selection algorithm as is the stabilized selected jet and
perceive represent the shapeless and other stabilizing device. Depending on the
stabilization of candidate, okay. Um, the target construction inference
guarantee is so. What is theta as tilde in this context?

\Editorial{2:19--2:45 is Woody reading slide/notation aloud; ASR is heavily
garbled. Retained as spoken. Sense: Woody asks Nick to define the notation,
especially \(\theta\) / \(\tilde{s}\), for the proposed coverage guarantee.}

\Omit{3:31--3:34 brief ask-to-repeat / connection remark while clarifying notation}

\Turn{3:37}{Nick Mino}
Oh, that's the probability that the selected object falls into the confidence
interval. So that it ends up being valid.

\Turn{3:57}{Woody Zhu}
Okay, so theta is basically the p-value, and you are hoping the p-value. Is
larger than or is it larger than one minus or like 0.95? Is that what you're
suggesting? I don't know. Yeah, Nick, can you define theta s and also the
confidence interval for us? I think there are some missing definitions here.

\Turn{4:25}{Nick Mino}
Yeah, theta s is the confidence interval. Here is saying that the interval where
s tilde is like correct, so the probability that theta tilde is in the
confidence interval should be hopefully one minus alpha, which would be like
95\%. So no, I mean, my question is: what is theta here?

\Editorial{Trailing ``So no, I mean, my question is: what is theta here?''
appears to be Woody in the source flow; retained as transcribed.}

\Turn{4:56}{Woody Zhu}
Because there's no formal. Explaining what theta is?

\Turn{5:06}{Nick Mino}
I don't remember specifically why I ended up using theta as the notation. I
think it was.

\Turn{5:12}{Woody Zhu}
Yeah, it's fine. I think overall, this direction is interesting. There are.
Large valued literature that has done something similar like this before. But I
think this idea in part aligns with one of the ideas Wenbin and I we have
discussed for quite a while. So I'm not sure if you still remember one of the
conversations that we have with Ryan regarding how do we define the stability
for a policy that is induced by optimization model. Yeah, I think that's what
exactly Nick wants to do. Okay. Yeah. But we probably wanted to make it more,
we have to have more clarity on the problem setup. I think this is a good
starting point, but we have to define everything we need more formally. In
terms of the notations and the problem setup. So, for example, you have defined
what the policy is, right? In this context, the not just like a random
data-dependent selection algorithm, it is a solution to an optimization
problem, right? So, if you remember that you, I think the policy is sort of
like takes as input some context variable. And outputs and prescribed decisions
by solving an optimization problem. And we have to, I think, one of the top
priorities is to formally define the stability of the policy induced by an
optimization problem. And once we have a clear definition of the stability for
this particular type of policies, then how do we find this stabilization
mechanism that will enable a valid optimization-based policy, which is be the
end of the objective of the study? I guess the end, one of the applications I
can think of is you can ask. Any kind of black box models or like AI systems
like ChatGPT to tell you what is the prediction of certain quantities. For
example, I wanted to know what is the stock price for Apple or for, I don't
know, a SpaceX tomorrow. You can certainly ask the same question to ChatGPT,
right? And the ChatGPT is give you some random number, or GPT can do some
reasoning and then give you some kind of more reasonable number that can be
used as the input to a downstream optimization model. And I would assume the
output of a GPT. Is supposed to be data dependent, right? Because the GPT was
trained on a specific large set of data. And we wanted to recalibrate the
output of this kind of combination of AI systems and optimization models so
that the output of this system is trustworthy. And one of the paper you can
look at. Is something actually I shared in the chat last night. You can search
it right now and Google is titled by no, it's called the prediction-powered
inference, right? Written by Angelopoulos and Michael Jordan. So, the basic
idea is pretty similar. So, they're trying to combine a powerful prediction
model. And some testing techniques in the way that the output can be
trustworthy. So, they wanted to assess the hypothesis with the help of an AI
model. So, that's basically what they do in this paper.

\Turn{9:34}{Woody Zhu}
Yeah, I think that's pretty much everything I can share, at least for now,
based on what I have written. So, two suggestions. One is you wanted to read
the paper first to get some high-level idea on the what, well, this is one of
the paper. The other paper is the, I assume you probably have read them before,
which is the algorithm of stability, right? Written by also Angelopoulos and
Michael Jordan.

\Turn{9:59}{Nick Mino}
Yeah, I think I've got it.

\Turn{10:01}{Woody Zhu}
Yeah, I think these are the two major papers that you probably wanted to read.
Yes, post-selection, right? Post- post-selection inference stability. So, you
probably want a closer look at both of these papers. And the top priority is to
set up the problem more formally. So, we will get a sense what kind of problem
we wanted to solve and what major technical challenges. And then we can come
back together again to have another meeting to discuss our next step.

\Turn{10:37}{Woody Zhu}
And in the meantime, you could think about maybe you can have a separate
discussion with Wenbin on possible applications, right? So where you need both
the AI system, like a ChatGPT, and optimization models in prescribing, for
example, high-stake decisions. I would assume, by the end of the day, we wanted
to. To set up the problem as a hypothesis testing.

\Turn{11:08}{Woody Zhu}
You can let the AI to prescribe a decision by producing some kind of random
quantity and then soft optimization, solve-by-authorization model. And then we
wanted to answer the question: should I trust this decision or not? So I think
that's going to be the end product I could envision right now.

\Editorial{``soft optimization, solve-by-authorization model'' is ASR garble;
sense in context: solve an optimization model / optimization-based decision.}

\Turn{11:35}{Woody Zhu}
Yeah, free to reach out to Wenbin if any, if you have any ideas or anything you
would like to discuss, once more concrete update or progress, just pin me in
the Slack chat so we can arrange another time to continue our discussion.

\Omit{11:54--12:34 assignment questions / SURA signature paperwork}

\Turn{12:35}{Nick Mino}
Oh, just like my current approach to trying to find the generalization has
been: here's like one of the compilations of the different. I have another one
that was here.

\Turn{12:58}{Nick Mino}
This is, yeah. So, like, I found as many like clean like formulas that then
were, or algorithms that then stabilized like selection.

\Turn{13:14}{Woody Zhu}
And I got a bunch of examples. Yeah, this is great. This is super helpful.

\Editorial{``And I got a bunch of examples.'' is likely Nick in the source flow;
ASR/source labeling is inconsistent here. Sense: Nick presented examples;
Woody praised them.}

\Turn{13:19}{Woody Zhu}
Maybe you can just share this document with us by email\ldots\ Yeah. Yeah, I
think, yeah, that's great.

\Omit{13:19 ASR fragment about adding Wenbin to the shared document}

\Turn{13:32}{Nick Mino}
I think this is a really good literature review and which can serve as a great
starting point.

\Editorial{13:32 is labeled Nick in the source. Content sounds like praise of
Nick's survey and may be Woody misattributed by ASR; left as in the source.}

\Turn{13:43}{Woody Zhu}
I think we have some discussion about it before, right? So the challenge of of
analyzing the stability for this type of policy lies at the discontinuity
induced by the optimization structure, if I understand correctly. So, because
if you, in the general setup, if I perturb the input, and correspondingly,
you're going to see the perturbation in the output space as well, right? So, I
think the perturbation should remain consistent and continuous. Across input
and output space. Whereas in the optimization problem, if you perturb the
input, it doesn't necessarily induce a perturbation in the output space because
the discrete nature of the decision boundary imposed by the optimization
structure itself.

\Turn{14:43}{Woody Zhu}
Yeah, so that's kind of the high-level intuition we have built up right now. We
know this problem is going to be non-trivial, but we also intuition how we want
to tackle this challenge because of the, one of the papers that Wenbin has done
before. So we sort of have some height of ideas, but we probably wanted into
the practice and give it more thoughts before we actually come up with the idea
how we want to solve this problem. But I think this is really good. This is a
really nice survey of the literature that has at least for some of the paper I
have read before and some of them I haven't read at all. Have you read all of
this paper or is just like a summarization done by AI?

\Editorial{``height of ideas'' / ``wanted into the practice'' left close to ASR;
sense: high-level ideas; dig into practice.}

\Turn{15:27}{Nick Mino}
No, I like looked for a bunch of papers and like tried to find examples and
then I also had AI do it. I spun up a bunch of agents and like to search through
a bunch to add more.

\Turn{15:45}{Nick Mino}
This has all of the citations for these. So you can go here and it'll bring you
to the link.

\Omit{15:39--15:44 fragmentary search for another slide/document}

\Turn{15:55}{Woody Zhu}
This is great. Yeah, this is great. This is fantastic.

\Turn{15:58}{Nick Mino}
And also, the little proofs on how they worked exactly.

\Turn{16:02}{Woody Zhu}
Oh, that's great. That's fantastic. Yeah. I think you're definitely on the very
good start point to work with. And just keep us posted if you have any
questions. And once you're ready, we can once we have a clear idea on the
problem setup. Maybe we can set up like a bi-weekly meeting with a fixed
schedule to catch up with the progress that you're making, while you're making
the progress. But at least for now, I think that's pretty much everything I can
think of.

\Omit{16:40--16:44 closing thanks}

\Takeaway{\textbf{CURRENT PLAN:} (1)~formalize stability for
optimization-induced policies; (2)~core challenge: discontinuity / discrete
decision boundaries under input perturbation; (3)~stabilization for valid
inference / trustworthy AI+optimizer pipelines; (4)~end product as hypothesis
test---trust the prescribed decision or not; (5)~immediate work: formal problem
setup; continue stabilized-selection literature survey; applications with
Wenbin; priority papers (prediction-powered inference; algorithmic /
post-selection stability).}

\clearpage
\section*{Appendix: Quick Reference Card}
\begin{center}
\begin{tabularx}{\linewidth}{@{}l >{\RaggedRight\arraybackslash}X@{}}
\toprule
\textbf{Item} & \textbf{Note} \\
\midrule
Extremely important sources & Meetings 1, 2, and 3 --- all three \\
Lab pillars & UQ / high-stakes decisions; grid resilience apps; human--AI collab \\
Chosen pillar & Algorithmic (CREDO, CREAM) \\
Superseded & Quantization engineering proposal \\
Superseded & Adaptive certified regret frontier (+ Overleaf formulation task) \\
Superseded (backup only) & Sequential CREAM; model/solver error accommodation \\
\textbf{Current plan} & Stability of optimization-induced policies; AI+opt trust /
  hypothesis testing \\
Key challenge & Optimization-induced discontinuity in decision boundaries \\
Priority papers (M3) & Prediction-powered inference; algorithmic / post-selection
  stability (Angelopoulos \& Jordan) \\
Background (M2, still useful) & Conformal prediction; LP / stochastic OR optimization \\
\bottomrule
\end{tabularx}
\end{center}

\section*{Archive Metadata}
\begin{tabular}{@{}ll@{}}
Document type & Edited verbatim research meeting transcript packet \\
Institution & Carnegie Mellon University \\
Student & Nicholas Mino \\
Faculty mentor & Woody Zhu \\
Graduate mentor & Wenbin Zhou \\
Meetings covered & 3 \\
Date range & 27 Apr 2026 -- 29 Jun 2026 \\
Compilation date & 26 July 2026 \\
\end{tabular}

\end{document}
